
% ============================================================
% PART III
% 3.2.1.2 -- GENERAL PROPERTIES
% BN-100 LUX
% ============================================================

\hypertarget{p3-2-1-2}{}

\subsubsection{General Properties}

% ============================================================
% 3.2.1.2.1 General Description
% ============================================================

\paragraph{General Description}

The BN-100 LUX is an active electrical neonatal radiant warmer
designed to provide thermal support to newborns and neonates.

The device uses non-contact infrared radiant heating. The heating
source is positioned above the infant bed, and a reflector assembly
directs radiant energy toward the infant. The control system
supports management of the heating function and monitoring of
skin temperature when the corresponding sensor and operating mode
are used.

The detailed operating characteristics and limitations are
described in the following subsections.

% ============================================================
% 3.2.1.2.2 Heating and Temperature Monitoring
% ============================================================

\paragraph{Heating and Temperature Monitoring}

According to the manufacturer documentation, the principal
heating and temperature-monitoring characteristics include:

\begin{itemize}
    \item heating tube rated power: \textbf{800 W};
    \item skin temperature control range:
    \textbf{$32$--$38\,^{\circ}\mathrm{C}$};
    \item stated overheating level 1:
    \textbf{$39\,^{\circ}\mathrm{C}$};
    \item stated overheating level 2:
    \textbf{$40\,^{\circ}\mathrm{C}$};
    \item stated temperature deviation alarm:
    \textbf{$\pm 1\,^{\circ}\mathrm{C}$};
    \item stated skin temperature sensor precision:
    \textbf{$\leq 0.3\,^{\circ}\mathrm{C}$}.
\end{itemize}

The heating mechanism is non-contact. Thermal energy is transferred
from the heating source toward the infant primarily through
radiant heat transfer.

The skin temperature sensor provides temperature information to
the control system and supports temperature monitoring and
automatic thermal control when the corresponding operating mode
is selected.

Correct sensor placement and attachment are important for reliable
temperature monitoring. The sensor shall be used in accordance
with the manufacturer's instructions for use, including the
specified precautions concerning incorrect placement or
detachment.

The stated temperature thresholds, deviation alarm and sensor
precision shall be interpreted together with the applicable
manufacturer specifications and technical documentation.

% ============================================================
% 3.2.1.2.3 Control and Alarm System
% ============================================================

\paragraph{Control and Alarm System}

The device incorporates electronic control circuitry associated
with heating, temperature monitoring and alarm functions.

The available manufacturer documentation identifies control
circuits, low-voltage circuits, relay and alarm circuitry, and
heating-control switching components. Detailed circuit diagrams,
component specifications, electrical tolerances and safety-related
circuit analyses are not reproduced in this subsection.

According to the manufacturer documentation, the specified
alarm conditions include:

\begin{itemize}
    \item mains power failure;
    \item skin temperature sensor failure;
    \item overheating; and
    \item temperature deviation.
\end{itemize}

The documentation also identifies an internal battery associated
with the alarm function during mains power interruption. The
battery configuration and the operation of the alarm during a
power failure are described in the applicable manufacturer
documentation.

Alarm functions, operating modes and associated precautions
shall be interpreted in accordance with the instructions for use
and the relevant technical documentation.

% ============================================================
% 3.2.1.2.4 Electrical Characteristics and Safety Classification
% ============================================================

\paragraph{Electrical Characteristics and Safety Classification}

According to the manufacturer documentation, the principal
electrical characteristics are:

\begin{itemize}
    \item rated supply voltage:
    \textbf{220 V $\pm 22$ V AC};
    \item rated frequency:
    \textbf{50 Hz $\pm 1$ Hz};
    \item maximum input power: \textbf{1000 VA}.
\end{itemize}

The manufacturer documentation identifies the following
electrical safety and operating characteristics:

\begin{itemize}
    \item protection against electric shock: \textbf{Class I};
    \item applied part classification: \textbf{Type BF};
    \item ingress protection rating: \textbf{IPX0};
    \item mode of operation: \textbf{Continuous operation}.
\end{itemize}

The manufacturer also states that the device is not intended for
use in the presence of flammable anaesthetic mixtures with air,
oxygen or nitrous oxide.

These characteristics concern electrical safety and operating
conditions. They shall not be confused with the regulatory
classification applicable to the Algerian homologation process.

Applicable electrical safety requirements and conformity
documentation are addressed in the relevant technical sections
of the dossier.

% ============================================================
% 3.2.1.2.5 Environmental Operating Conditions
% ============================================================

\paragraph{Environmental Operating Conditions}

The specified environmental conditions for operation are:

\begin{itemize}
    \item \textbf{Ambient temperature:}
    $+18\,^{\circ}\mathrm{C}$ to $+30\,^{\circ}\mathrm{C}$;
    \item \textbf{Relative humidity:} 30\% to 75\%;
    \item \textbf{Atmospheric pressure:} 700 to 1060 hPa;
    \item \textbf{Air velocity:} less than 0.3 m/s.
\end{itemize}

The device should be operated within the specified environmental
limits to maintain its intended performance and temperature
measurement characteristics.

% ============================================================
% 3.2.1.2.6 Storage and Transportation Conditions
% ============================================================

\paragraph{Storage and Transportation Conditions}

The specified storage and transportation conditions are:

\begin{itemize}
    \item \textbf{Temperature:}
    $-40\,^{\circ}\mathrm{C}$ to $+55\,^{\circ}\mathrm{C}$;
    \item \textbf{Relative humidity:} not exceeding 95\%;
    \item \textbf{Atmospheric pressure:} 500 to 1060 hPa.
\end{itemize}

The device should be protected against environmental conditions
that could adversely affect its mechanical, electrical or
electronic components during storage and transportation.

% ============================================================
% 3.2.1.2.7 Mechanical Characteristics and Load Limits
% ============================================================

\paragraph{Mechanical Characteristics and Load Limits}

The BN-100 LUX incorporates an infant bed and auxiliary tray
structures with specified load limits.

According to the manufacturer documentation:

\begin{itemize}
    \item maximum infant bed load: \textbf{100 N},
    corresponding approximately to \textbf{10 kg};
    \item maximum tray load: \textbf{20 N},
    corresponding approximately to \textbf{2 kg}.
\end{itemize}

These limits apply to their respective structural components
and shall not be interpreted as interchangeable load ratings.

% ============================================================
% 3.2.1.2.8 Service Life
% ============================================================

\paragraph{Service Life}

The specified expected service life of the BN-100 LUX is
\textbf{8 years}. The specified service life of the heating
tube is \textbf{25,000 hours}.

These service-life specifications do not eliminate the need
for preventive maintenance, inspection, cleaning, functional
verification and component replacement where required.

Actual operational life depends on the conditions of use,
maintenance practices and the condition of safety-critical
components.

% ============================================================
% 3.2.1.2.9 Reusability and Cleaning
% ============================================================

\paragraph{Reusability and Cleaning}

The BN-100 LUX is a reusable medical device and is not intended
to be supplied as a sterile, single-use medical device.

Cleaning and disinfection shall be performed in accordance
with the manufacturer's instructions for use and applicable
institutional procedures.

Cleaning and disinfection agents shall be selected with due
consideration for their compatibility with the device materials
and surfaces.

% ============================================================
% 3.2.1.2.10 Patient Safety-Related Characteristics
% ============================================================

\paragraph{Patient Safety-Related Characteristics}

The device incorporates operational precautions intended to
reduce risks associated with neonatal radiant heating.

The manufacturer's safety information includes:

\begin{itemize}
    \item warnings concerning excessive heating and potential burns;
    \item precautions related to incorrect placement or detachment
    of the skin temperature sensor;
    \item precautions concerning oxygen-enriched environments
    and fire hazards;
    \item precautions associated with manual heating modes; and
    \item precautions associated with optional phototherapy
    functions, including protection of the infant's eyes.
\end{itemize}

These precautions shall be considered together with the
manufacturer's instructions for use and risk-management
documentation.

% ============================================================
% 3.2.1.2.11 Accessories and Optional Modules
% ============================================================

\paragraph{Accessories and Optional Modules}

Depending on the supplied configuration, the BN-100 LUX may
include the following accessories or optional modules:

\begin{itemize}
    \item LED phototherapy unit;
    \item spotlight;
    \item X-ray cassette tray;
    \item intravenous support;
    \item oxygen cylinder support;
    \item skin temperature sensor; and
    \item low-pressure suction equipment.
\end{itemize}

The applicable configuration and included accessories are
identified in the corresponding product and packing
documentation.

% ============================================================
% 3.2.1.2.12 Material Information
% ============================================================

\paragraph{Material Information}

The manufacturer documentation identifies the following
materials for specific components:

\begin{itemize}
    \item carbon steel for the radiant heating tube;
    \item Lucite/Plexiglas for the side baffles; and
    \item polycarbonate for the suction bottle.
\end{itemize}

These material identifications apply to the components
specified in the manufacturer documentation and do not
constitute a complete material declaration for the entire
device.

% ============================================================
% 3.2.1.2.13 Applicable Standards and Verification Evidence
% ============================================================

\paragraph{Applicable Standards and Verification Evidence}

The manufacturer documentation identifies the following
standards and technical requirements in relation to the device:

\begin{itemize}
    \item IEC 60601-1;
    \item IEC 60601-1-8;
    \item IEC 60601-2-21;
    \item CISPR 11, Group 1, Class A;
    \item IEC 61000-4-2;
    \item IEC 61000-4-3;
    \item IEC 61000-4-4; and
    \item IEC 61000-4-5.
\end{itemize}

These references provide the technical framework for the
relevant safety, alarm, electromagnetic compatibility and
radiant-warmer requirements. The associated conformity
documentation and test evidence are addressed in the relevant
technical sections of the dossier.

% ============================================================
% END OF 3.2.1.2
% ============================================================
